\section{Assemblées générales}
\label{assemblees-generales}

\subsection{Lignes directrices des réunions}
\label{assemblees-generales-lignes-directrices-des-reunions}

\subsubsection{Convocation d'une assemblée}
\label{lignes-directrices-des-reunions-convocation-d-une-assemblee}

\begin{enumerate}
 \item Une assemblée générale sera convoquée sur résolution du CA ou sur demande écrite d'au moins 5 \% des membres.
\end{enumerate}

\subsubsection{Traduction en direct}
\label{lignes-directrices-des-reunions-traduction-en-direct}

\begin{enumerate}
 \item À chaque assemblée générale, des services de traduction simultanée seront disponibles, qu'ils soient confiés à des ressources professionnelles externes ou dirigés par les membres.
\end{enumerate}

\subsubsection{Avis de convocation}
\label{lignes-directrices-des-reunions-avis-de-convocation}

\begin{enumerate}
 \item Un avis écrit de vingt-et-un (21) jours sera donné aux membres pour toute assemblée générale.
 \item Aucune erreur ou omission dans la transmission de l'avis de toute assemblée générale de la FCEG n'invalidera une telle assemblée ni n'annulera les procédures qui y seront prises.
\end{enumerate}

\subsubsection{Vote et quorum}
\label{lignes-directrices-des-reunions-vote-et-quorum}

\begin{enumerate}
 \item Tous les membres auront un (1) vote.
 \item Toutes les questions soumises à l'assemblée générale seront décidées par une majorité des votes des membres présents et votants, sauf indication contraire.
 \item Les deux tiers (2/3) des membres constitueront le quorum lors d'une assemblée générale.
\end{enumerate}

\subsubsection{Droit de parole}
\label{lignes-directrices-des-reunions-droit-de-parole}

\begin{enumerate}
 \item Les membres et l'équipe de direction auront le droit de parole lors des assemblées générales.
\end{enumerate}

\subsubsection{Ordre}
\label{lignes-directrices-des-reunions-ordre}

\begin{enumerate}
 \item Lors de toutes les assemblées générales, la présidence du CA doit appliquer les règles des assemblées délibérantes telles que décrites dans les Robert’s Rules of Order, sauf lorsque des dispositions contraires sont adoptées par la FCEG.
 \item La présidence du CA présidera les assemblées générales.
\end{enumerate}

\subsubsection{Personne représentant}
\label{lignes-directrices-des-reunions-personne-representant}

\begin{enumerate}
 \item Un membre doit, au moyen d'une communication vérifiable par sa présidence, nommer une personne représentante pour assister et agir lors d'une assemblée générale spécifique.
 \item L'avis de chaque assemblée générale doit rappeler aux membres leur droit de nommer une personne représentante.
 \item Toutes les personnes représentantes doivent être une personne étudiante membre de l'association membre ou d’une autre association membre de la FCEG.
 \item Toutes les réglementations relatives au vote par procuration doivent également se conformer à l'article 74 du Règlement sur les organisations à but non lucratif de régime fédéral.
\end{enumerate}

\subsubsection{Procès-verbaux}
\label{lignes-directrices-des-reunions-proces-verbaux}

\begin{enumerate}
 \item Des procès-verbaux seront rédigés en anglais et en français lors de toutes les assemblées générales.
 \item Le texte complet et exact des procès-verbaux sera distribué à tous les membres en anglais et en français dans les trente jours suivant la levée de l'assemblée.
\end{enumerate}

\subsection{Vidéos précédant l'assemblée générale}
\label{assemblees-generales-videos-precedant-l-assemblee-generale}

\begin{enumerate}
 \item L'exécutif national peut créer des vidéos avant toute activité impliquant une assemblée générale afin de s'assurer que les membres sont informés des sujets pertinents pour le congrès, y compris, mais sans s'y limiter, les politiques en cours de discussion et les mandats en cours.
 \item La durée totale de ces vidéos ne dépassera pas 30 minutes, et les vidéos doivent être fournies au moins 1 semaine avant le début du congrès.
\end{enumerate}
